\documentclass[10pt]{article}
\usepackage[T1]{fontenc}
\usepackage[utf8]{inputenc}
\usepackage[french]{babel}
\usepackage{amsmath,amssymb,mathrsfs,mathpazo}
\usepackage{geometry}
\geometry{a4paper,margin=2cm}
\usepackage{graphicx,array,tabularx,tabularray,colortbl,multirow,multicol,booktabs,enumitem,pifont,ulem,fancybox,diagbox,cancel}
\usepackage{pstricks,pstricks-add,pst-plot,pst-tree,pst-node,pst-text,pst-func,pst-math,pst-bspline,pst-3dplot,pst-eucl,pst-solides3d}
\let\pscalculate\undefined
\usepackage{pst-bezier}
\usepackage{tikz,pgfplots,tkz-tab}
\usetikzlibrary{arrows,arrows.meta,patterns,calc,positioning,trees,decorations.markings,shapes.geometric,shapes.misc,shapes.arrows,intersections,decorations.pathmorphing,decorations.pathreplacing,backgrounds,shadows,babel}
\pgfplotsset{compat=1.18}
\usepgfplotslibrary{fillbetween}
\usepackage[np]{numprint}
\usepackage{listings,makecell,pas-tableur,twoopt}
\let\barre\undefined
\usepackage[pstricks]{tablvar}
\definecolor{darkgreen}{RGB}{0,100,0}
\providecommand{\red}{\color{red}}
\providecommand{\blue}{\color{blue}}
\providecommand{\green}{\color{green}}
\providecommand{\black}{\color{black}}
\providecommand{\white}{\color{white}}
\providecommand{\gray}{\color{gray}}
\providecommand{\cyan}{\color{cyan}}
\providecommand{\magenta}{\color{magenta}}
\providecommand{\yellow}{\color{yellow}}
\providecommand{\violet}{\color{violet}}
\providecommand{\orange}{\color{orange}}
\providecommand{\euro}{\texteuro}
\providecommand{\EUR}{\texteuro}
\providecommand{\ieme}{\textsuperscript{e}}
\providecommand{\degres}{\ensuremath{^\circ}}
\providecommand{\degree}{\ensuremath{^\circ}}
\providecommand{\vv}[1]{\overrightarrow{#1}}
\providecommand{\Coord}[1]{#1}
\providecommand{\decosix}{}
\providecommand{\starredbullet}{\ensuremath{\bullet}}
\pagestyle{empty}
\setlength{\parindent}{0pt}
\usepackage{adjustbox}\DeclareUnicodeCharacter{2236}{:}\DeclareUnicodeCharacter{2190}{\ensuremath{\leftarrow}}
\providecommand{\R}{\mathbb{R}}
\providecommand{\N}{\mathbb{N}}
\providecommand{\Z}{\mathbb{Z}}
\providecommand{\Q}{\mathbb{Q}}
\providecommand{\D}{\mathbb{D}}
\providecommand{\C}{\mathbb{C}}
\providecommand{\e}{{\mathrm{e}}}
\providecommand{\vect}[1]{\overrightarrow{#1}}
\providecommand{\vectt}[1]{\overrightarrow{\mathrm{#1}}}
\providecommand{\vvt}[1]{\overrightarrow{\mathrm{#1}}}
\providecommand{\np}[2][]{{#2~#1}}
\providecommand{\vv}[1]{{\overrightarrow{#1}}}
\renewcommand{\d}{{\mathrm{d}}}
\providecommand{\degres}{{^\circ}}
\providecommand{\degre}{{^\circ}}
\providecommand{\blue}{\color{blue}}
\providecommand{\red}{\color{red}}
\providecommand{\enskip}{\quad}
\providecommand{\llbracket}{\lbrack\!\lbrack}
\providecommand{\rrbracket}{\rbrack\!\rbrack}
\providecommand{\euro}{{\text{€}}}
\providecommand{\no}{{\text{n°}}}
\providecommand{\interval}[2]{{[#1;#2]}}
\providecommand{\intervFO}[2]{{[#1;#2[}}
\providecommand{\dotfill}{\ldots}

\providecommand{\ds}{\displaystyle}
\providecommand{\textding}[1]{#1}
\providecommand{\Oij}{\(\mathrm{O};\vec{\imath},\vec{\jmath}\)}
\providecommand{\Oijk}{\(\mathrm{O};\vec{\imath},\vec{\jmath},\vec{k}\)}
\providecommand{\Ouv}{\(\mathrm{O};\vec{u},\vec{v}\)}
\providecommand{\intervFF}[2]{\ensuremath{[#1;#2]}}
\renewcommand{\euro}{\ifmmode\text{\texteuro}\else\texteuro\fi}
\DeclareUnicodeCharacter{00B0}{\ifmmode^{\circ}\else\textdegree\fi}
\DeclareUnicodeCharacter{20AC}{\ifmmode\text{\texteuro}\else\texteuro\fi}
% Énoncés : Éducation nationale. Transcription et corrigés : APMEP, Denis Vergès et contributeurs. Découpage et classement : CoopMaths, Sébastien Lozano, Éric Elter et contributeurs. Adaptation de la présentation : Maths Explorer.
% CC BY-SA : https://coopmaths.fr/www/about/
% https://coopmaths.fr/alea/static/bac/2022/tex/bac_2022_05_sujet2_asie_3_cor.tex
\begin{document}
\noindent\textit{APMEP / CoopMaths. CC BY-SA. Présentation adaptée par Maths Explorer.}\par\medskip
\medskip

\emph{Principaux domaines abordés}: Probabilités conditionnelles et indépendance. Variables aléatoires.

\bigskip

\begin{center}\emph{Les deux parties de cet exercice sont indépendantes.}\end{center}

\textbf{Partie 1}

\medskip
























\begin{enumerate}
\item 
S'il a pris le bus il est à l'heure pour son vol, donc $P_B(V) = 1$.

\item On représente la situation par un arbre pondéré.

\begin{center}
\pstree[treemode=R,nodesepB=3pt,levelsep=3cm]{\TR{}}
{\pstree{\TR{$B$~~} \taput{$0,2$}}
	{\TR{$V$}\taput{1}
	\TR{$\overline{V}$}\tbput{0}
	}
\pstree{\TR{$\overline{B}$~~} \tbput{$0,8$}}
	{\TR{$V$}\taput{0,5}
	\TR{$\overline{V}$}\tbput{0,5}
	}	
}
\end{center}

\item 
D'après la loi des probabilités totales :

$P(V) = P(B \cap V) + P\left(\overline{B} \cap V\right) = 0,2 \times 1 + 0,8 \times 0,5 = 0,2 + 0,4 = 0,6$.

\item 
On calcule $P_V(B) = \dfrac{P(V \cap B)}{P(V)} = \dfrac{P(B \cap V)}{P(V)} = \dfrac{0,2}{0,6} = \dfrac13$.
\end{enumerate}

\bigskip

\textbf{Partie 2}

\medskip











\begin{enumerate}
\item 
La présence à l'embarquement de chaque passager est indépendante des autres et chaque passager a la même probabilité 0,95 d'être présent, donc la variable aléatoire $X$ suit la loi binomiale de paramètres $n = 206$ et $p = 0,95$.
\item 
On a E $ = n \times p = 206 \times 0,95 = 195,7 \approx 196$.

En moyenne sur 206 titulaires d'un billet à peu près 196 vont se présenter.
\item 
On a $P(X = 201) = \binom{206}{201} \times 0,95^{201} \times 0,05^{5} \approx \np{0,03063}$, soit 0,031 au millième près.
\item 
La calculatrice donne $P(X \leqslant 200) \approx \np{0,9477}$, soit 0,948 au millième près.

 Il est donc à peu près certain que l'avion sera au mieux juste rempli.
\item 












On admet que $Y$ suit la loi de probabilité donnée par le tableau suivant:

\begin{center}
\begin{tabularx}{\linewidth}{|c|*{7}{>{\centering \arraybackslash}X|}}\hline
$y_i$					&0	&1	&2	&3	&4	&5	&6\\ \hline
$P\left(Y = y_i\right)$	&\np{0,94775}&\np{0,03063}&\np{0,01441}&\np{0,00539}&\np{0,00151}&\np{0,00028}&$\blue\np{0,00003}$\\ \hline
$c_i$&$\blue\np{51500}$&$\blue\np{50650}$&$\blue\np{49800}$&$\blue\np{48950}$&$\blue\np{48100}$&$\blue\np{47250}$&$\blue\np{46400}$\\ \hline
\end{tabularx}
\end{center}

	\begin{enumerate}
		\item 
En complétant à 1 la somme des probabilités données dans le tableau on trouve:\\
 $P(Y = 6) =1 - \left ( P(Y=0) + P(Y=1) + \cdots + P(Y=5)\strut\right )= \np{0,00003}$.
 
		\item 
La compagnie a encaissé $206 \times 250 = \np{51500}$~\euro{} et elle devra rembourser 850~\euro{} à chaque client ne pouvant embarquer, donc $C = \np{51500} - 850Y$
		\item 
Voir le tableau ci-dessus.


On a E$(C) = \np{51500} \times \np{0,94775} + \ldots \np{46400} \times \np{0,00003} \approx \np{51429,2}$, soit \np{51429}~\euro{} à l'euro près

		\item 
En vendant exactement 200 billets, la compagnie fera un chiffre d'affaires de $200\times 250$ soit $\np{50000}$ euros.

En pratiquant le surbooking, la compagnie peut espérer un chiffre d'affaires de \np{51429} euros.
	\end{enumerate}
\end{enumerate}

\bigskip

\end{document}
